\SourceLeaf{109}%
nen Resultaten eine grössere Genauigkeit beilegt, als sie wirklich besitzen. Freilich hat in dem Falle, wo die Anzahl der Beobachtungen vielemale grösser ist als die der unbekannten Grössen, diese Unrichtigkeit wenig zu bedeuten; allein theils erfordert die Würde der Wissenschaft, dass man vollständig und bestimmt übersehe, wieviel man hierdurch zu fehlen Gefahr läuft, theils sind auch wirklich öfters nach jenem fehlerhaften Verfahren Rechnungsresultate in wichtigen Fällen aufgestellt, wo jene Voraussetzung nicht Statt fand. Der Verfasser hat daher diesen Gegenstand einer besondern Untersuchung unterworfen, die zu einem sehr merkwürdigen höchst einfachen Resultate geführt hat. Man braucht nemlich den nach dem angezeigten fehlerhaften Verfahren gefundenen mittlern Fehler, um ihn in den richtigen zu verwandeln, nur mit
\[
\surd\frac{\pi-\rho}{\pi}
\]
zu multipliciren, wo \(\pi\) die Anzahl der Beobachtungen und \(\rho\) die Anzahl der unbekannten Grössen bedeutet.

Die letzte Untersuchung betrifft noch die Ausmittelung des Grades von Genauigkeit, welcher dieser Bestimmung des mittlern Fehlers selbst beigelegt werden muss: die Resultate derselben müssen aber in der Abhandlung selbst nachgelesen werden.

